% =============================================================================
\section{Motivation: The Limits of Static Word Embeddings}
\label{sec:static}
% =============================================================================

\subsection{What Are Word Embeddings?}

The first major breakthrough beyond symbolic NLP was the introduction of
\emph{distributed word representations}, or \emph{word embeddings}: mappings
from vocabulary items to dense vectors in a continuous $d$-dimensional space
$\R^d$.

The foundational idea is the \textbf{distributional hypothesis}~\cite{harris1954distributional}:
words that occur in similar contexts tend to carry similar meanings. Models such
as Word2Vec~\cite{mikolov2013word2vec}, GloVe~\cite{pennington2014glove}, and
FastText~\cite{bojanowski2017fasttext} operationalise this by training on large
text corpora to produce an embedding matrix $\mathbf{E} \in \R^{|V| \times d}$,
where each row $\mathbf{e}_w$ is the embedding of word $w$.

\textbf{Word2Vec}~\cite{mikolov2013word2vec} (Google, 2013) trains a shallow
neural network with either a Skip-gram or CBOW (Continuous Bag-of-Words)
objective to predict a word from its context or vice versa.
\textbf{GloVe}~\cite{pennington2014glove} (Stanford, 2014) learns embeddings
from the global word co-occurrence matrix.
\textbf{FastText}~\cite{bojanowski2017fasttext} (Facebook, 2016) extends
Word2Vec to subword units, improving handling of morphologically rich languages
and out-of-vocabulary words.

\begin{table}[H]
\centering
\caption{Summary of static word embedding models.}
\label{tab:static}
\begin{tabularx}{0.92\textwidth}{>{\bfseries}lllX}
  \toprule
  Model & Origin & Year & Key Innovation \\
  \midrule
  Word2Vec~\cite{mikolov2013word2vec}    & Google   & 2013
    & Skip-gram / CBOW objectives; scalable training \\
  GloVe~\cite{pennington2014glove}       & Stanford & 2014
    & Global co-occurrence matrix factorisation \\
  FastText~\cite{bojanowski2017fasttext} & Facebook & 2016
    & Subword $n$-gram representations \\
  \bottomrule
\end{tabularx}
\end{table}

\noindent
These models yield rich geometric properties. The famous analogy
$\vec{\text{king}} - \vec{\text{man}} + \vec{\text{woman}} \approx
\vec{\text{queen}}$ emerges naturally from the learned space, illustrating that
semantic and syntactic relationships are captured as linear offsets in the
embedding space.

\subsection{The Fatal Flaw: Context-Blindness}

Despite their utility, static embeddings have one fundamental limitation:
\textbf{every occurrence of a word receives the same vector}, regardless of the
sentence it appears in. Formally, the mapping $f\colon w \mapsto \mathbf{e}_w$
is a simple lookup --- context plays no role.

For the polysemy example above, both occurrences of \textit{bank} map to an
identical vector $\mathbf{e}_{\text{bank}} \in \R^d$. The embedding must
somehow encode both the financial-institution sense and the geographical sense
simultaneously, which it cannot do faithfully. This context-blindness makes
static embeddings inadequate for tasks that require genuine language
understanding.
